\chapter{Relation Between Lattice And Continuum Local QFT}
\label{ch:constructive-lattice-continuum-local-qft}

\ConventionsEuclidean

This chapter explains how finite-lattice data become continuum local QFT
data.  The passage requires scaling maps for observables, convergence of
correlation functions as distributions, reflection positivity, and
reconstruction or local-net comparison.

\section{Scaling Maps}

Let \(a>0\) be the lattice spacing and let \(\Lambda_a=a\mathbb Z^d\cap V\)
inside a finite region \(V\).  A lattice field \(\phi_a(x)\) becomes a
continuum distribution only after a scaling map is specified:
\[
  \Phi_a(f)
  =
  Z_a^{1/2}a^d\sum_{x\in\Lambda_a} f(x)\phi_a(x),
\]
where \(f\in C_c^\infty(V)\) and \(Z_a\) is the field normalization.  Composite
operators require a mixing matrix
\[
  \mathcal O_{A,a}^{\mathrm{ren}}
  =
  \sum_B Z_{AB}(a)\mathcal O_{B,a}.
\]
The operator basis and normalization are part of the continuum-limit datum.

\section{Distributional Convergence}

The \(n\)-point Schwinger distributions are
\[
  S_{n,a}(f_1,\ldots,f_n)
  =
  \left\langle
  \Phi_a(f_1)\cdots \Phi_a(f_n)
  \right\rangle_a .
\]
A continuum Schwinger function \(S_n\) is obtained when
\[
  \lim_{a\to0}S_{n,a}(f_1,\ldots,f_n)
  =
  S_n(f_1,\ldots,f_n)
\]
for all test functions in the chosen class, with compatible infinite-volume
and continuum limits.  Pointwise lattice correlators are auxiliary data; the
continuum field is distributional.

\section{Reflection Positivity And OS Reconstruction}

Assume the lattice measures are reflection positive with respect to a
hyperplane and that the scaling limit preserves reflection positivity:
\[
  \sum_{i,j}\overline c_i c_j\,S(F_i^\Theta F_j)\geq0
\]
for all finite linear combinations of positive-time test functionals.  If the
limit Schwinger functions also satisfy Euclidean covariance, symmetry,
regularity, and the OS growth hypotheses, OS reconstruction produces a
Hilbert space, vacuum, Hamiltonian, and Wightman distributions.  Thus the
lattice-to-continuum bridge is a theorem only after these hypotheses are
verified for the limiting distributions.

\section{Local Algebras From Lattice Algebras}

For a bounded continuum region \(O\), let \(O_a\subset\Lambda_a\) be the set
of lattice sites or links whose cells lie in \(O\).  The lattice observable
algebra \(\mathcal A_a(O_a)\) is finite dimensional in spin systems and
finite-regulator gauge systems after gauge constraints are imposed.  A
continuum local algebra is approached by embeddings or correlation-function
limits
\[
  \mathcal A_a(O_a)\longrightarrow \mathcal A(O)
\]
that preserve inclusion, locality, and the state in the scaling limit.  The
construction is delicate for gauge theories because Gauss-law constraints and
center degrees of freedom affect the algebra assigned to a region with
boundary.

\section{Gauge-Invariant Operators}

In lattice gauge theory, link variables are regulator coordinates.  Continuum
local gauge-invariant operators arise from small Wilson loops, covariant
lattice differences, and matter bilinears after subtracting and mixing all
operators with the same symmetries.  For example, the continuum action
density is represented by a plaquette combination
\[
  \mathcal P_a(x)
  =
  \frac{1}{a^4}\left(1-\frac{1}{\dim R}\operatorname{Re}\operatorname{Tr}_R
  U_{\square_x}\right)
  +\text{local mixings}.
\]
The mixings include identity and lower-dimensional terms allowed by the
regulator symmetries.

\section{Universality As A Convergence Statement}

Two lattice regulator families define the same continuum QFT when there are
scaling maps for their observables such that all reconstructed local QFT data
agree: Schwinger distributions, Wightman distributions or local nets,
symmetry actions, stress tensor normalization, line or defect sectors when
included, and background responses.  Equality of critical exponents is a
useful diagnostic but is not the full equivalence relation.

\section{Continuum Limit Checklist}

A lattice construction intended to define a continuum QFT must record:
\[
\begin{array}{ll}
\text{regulator family} & \text{lattice actions, measures, boundary
conditions},\\
\text{scaling trajectory} & \text{couplings as functions of }a,\\
\text{operator maps} & \text{normalization and mixing},\\
\text{limit topology} & \text{distributional, OS, Wightman, or local-net},\\
\text{positivity} & \text{reflection positivity or replacement structure},\\
\text{symmetries} & \text{continuum symmetry recovery and anomalies}.
\end{array}
\]

\begin{openproblem}[Local nets from interacting lattice gauge limits]
\label{op:local-nets-from-lattice-gauge-limits}
Construct continuum local algebras for interacting lattice gauge theory
scaling limits, including Gauss-law center choices, charged sectors,
extended operators, reflection positivity, and comparison with Wightman or OS
reconstruction.
\end{openproblem}
